%============================================================
% macros.tex — short commands used throughout the text
%============================================================

% --- Links ---
% blue "Solve" button (\exbadge comes from preamble/difficulty.tex)
\newcommand{\solveit}[1]{\href{#1}{\exbadge{primaryblue}{Solve $\nearrow$}}}
\newcommand{\bloglink}[1]{\href{#1}{\textcolor{successgreen}{\texttt{[Blog]}}}}
\newcommand{\vidlink}[1]{\href{#1}{\textcolor{accentorange}{\texttt{[Video]}}}}

% --- DP notation ---
\newcommand{\dpstate}[1]{\texttt{\textcolor{deepblue}{dp[#1]}}}
\newcommand{\memo}[1]{\texttt{\textcolor{memoorange}{memo[#1]}}}
\newcommand{\bigO}[1]{\ensuremath{\mathcal{O}(#1)}}
\newcommand{\bigT}[1]{\ensuremath{\Theta(#1)}}
\newcommand{\bigOm}[1]{\ensuremath{\Omega(#1)}}
\newcommand{\INF}{\ensuremath{+\infty}}
\newcommand{\contradiction}{\ensuremath{\bot}}

% --- Graph notation ---
\newcommand{\nV}{\ensuremath{|V|}}
\newcommand{\nE}{\ensuremath{|E|}}
\newcommand{\VE}{\ensuremath{\mathcal{O}(|V|+|E|)}}
\newcommand{\dist}[1]{\ensuremath{\operatorname{dist}(#1)}}
\newcommand{\Adj}[1]{\ensuremath{\operatorname{adj}(#1)}}
\newcommand{\indeg}[1]{\ensuremath{\deg^{-}(#1)}}
\newcommand{\outdeg}[1]{\ensuremath{\deg^{+}(#1)}}
\newcommand{\tin}[1]{\ensuremath{\mathrm{tin}[#1]}}
\newcommand{\tout}[1]{\ensuremath{\mathrm{tout}[#1]}}
\newcommand{\low}[1]{\ensuremath{\mathrm{low}[#1]}}
\newcommand{\lca}[2]{\ensuremath{\operatorname{lca}(#1,#2)}}
\newcommand{\findset}[1]{\ensuremath{\operatorname{find}(#1)}}
% --- Array notation ---
\newcommand{\seg}[2]{\ensuremath{[#1, #2)}}            % half-open range [l, r)
\newcommand{\cseg}[2]{\ensuremath{[#1, #2]}}           % closed range [l, r]
\newcommand{\pre}[1]{\ensuremath{P[#1]}}               % prefix sum
\newcommand{\sumr}[2]{\ensuremath{\operatorname{sum}(#1, #2)}}
\newcommand{\mex}{\ensuremath{\operatorname{mex}}}
% cross-references to the companion books:  \dpbook{1}{Kadane}   \treebook{13}{Fenwick Trees}
\newcommand{\dpbook}[2]{\textit{Dynamic Programming} book, Ch.~#1 (#2)}
\newcommand{\graphbook}[2]{\textit{Graph Algorithms} book, Ch.~#1 (#2)}
\newcommand{\treebook}[2]{\textit{Trees} book, Ch.~#1 (#2)}
\newcommand{\stringbook}[2]{\textit{Strings} book, Ch.~#1 (#2)}
\newcommand{\recbook}[2]{\textit{Recursion} book, Ch.~#1 (#2)}
\newcommand{\arraybook}[2]{\textit{Arrays} book, Ch.~#1 (#2)}
\newcommand{\bitbook}[2]{\textit{Bit Manipulation} book, Ch.~#1 (#2)}
\newcommand{\sqhbook}[2]{\textit{Stacks, Queues \& Heaps} book, Ch.~#1 (#2)}

% --- Verdicts ---
\newcommand{\tle}{\textcolor{warningred}{\textbf{TLE}}}
\newcommand{\ac}{\textcolor{successgreen}{\textbf{AC}}}
\newcommand{\mle}{\textcolor{warningred}{\textbf{MLE}}}
\newcommand{\wa}{\textcolor{warningred}{\textbf{WA}}}
\newcommand{\stackoverflow}{\textcolor{warningred}{\textbf{Stack Overflow}}}

% --- Approach tags ---
\newcommand{\recursive}{\textcolor{recursionred}{\textbf{[Recursion]}}}
\newcommand{\memoized}{\textcolor{memoorange}{\textbf{[Memoization]}}}
\newcommand{\tabulated}{\textcolor{tabulationblue}{\textbf{[Tabulation]}}}
\newcommand{\optimized}{\textcolor{optimizedpurple}{\textbf{[Space-Optimized]}}}

% --- Difficulty / rating ---
% \stars{filled}{total}
\newcommand{\stars}[2]{%
    \foreach \i in {1,...,#1}{\textcolor{accentorange}{\ding{72}}}%
    \pgfmathtruncatemacro{\emptycount}{#2 - #1}%
    \ifnum\emptycount>0
        \foreach \i in {1,...,\emptycount}{\textcolor{gray!40}{\ding{72}}}%
    \fi}
\newcommand{\easy}{\textcolor{successgreen}{\textbf{Easy}}}
\newcommand{\medium}{\textcolor{accentorange}{\textbf{Medium}}}
\newcommand{\hard}{\textcolor{warningred}{\textbf{Hard}}}

% --- Check marks ---
\newcommand{\yes}{\textcolor{successgreen}{\ding{51}}}
\newcommand{\no}{\textcolor{warningred}{\ding{55}}}
\newcommand{\todo}[1]{\textcolor{warningred}{\textbf{[TODO: #1]}}}

% --- Small caption under a diagram (works inside boxes) ---
\newcommand{\figcap}[1]{%
    \par\smallskip{\centering\small\itshape\color{darkgray}#1\par}\smallskip}

% --- Section-level helpers inside boxes ---
\newcommand{\boxsection}[1]{\subsection*{#1}}
\newcommand{\boxsubsection}[1]{\subsubsection*{#1}}
\newcommand{\step}[1]{\textbf{\textcolor{deepblue}{#1}}}

% --- Theorem environments (outside boxes) ---
\theoremstyle{definition}
\newtheorem{lemma}{Lemma}[chapter]
\newtheorem{claim}{Claim}[chapter]

% --- Labels of solved problems ---
% \problemlabel{1.4.3}{ssec:ch1-stock-iv} replaces \label right after \subsection{Problem 1.4.3: ...}.
% A plain \label would make \ref print the subsection counter (e.g. 1.5.4), which is not the problem number.
\makeatletter
\newcommand{\problemlabel}[2]{\protected@edef\@currentlabel{#1}\label{#2}}
\makeatother
